\documentclass[10pt,journal,compsoc]{IEEEtran}

\usepackage{amsmath,amssymb,amsfonts}
\usepackage{algorithmic}
\usepackage{graphicx}
\usepackage{textcomp}
\usepackage{xcolor}
\usepackage{booktabs}
\usepackage{multirow}
\usepackage{microtype}
\usepackage{hyperref}

\hypersetup{
    colorlinks=true,
    linkcolor=blue,
    citecolor=blue,
    urlcolor=blue
}

\begin{document}

\title{SEFA-Forensics: Self-Eliciting Foundation Anchors for Multi-Dataset Inpainting and Splicing Forensics Without Backbone Fine-Tuning}

\author{Akash Biswas,~\IEEEmembership{Graduate Student Member,~IEEE}
\thanks{Manuscript received October 17, 2026. The author is with the Department of Computer Science and Engineering.}
}

\markboth{IEEE Transactions on Information Forensics and Security,~Vol.~XX, No.~X, 2026}%
{Biswas: SEFA-Forensics: Multi-Dataset Inpainting and Splicing Forensics}

\IEEEtitleabstractindextext{%
\begin{abstract}
Detecting subtle image inpainting, copy-move manipulation, and deep generative synthesis remains an open challenge due to cross-generator domain shifts and post-processing laundering. While large vision foundation models (such as DINOv2) encode rich natural image priors, conventional forensic fine-tuning either destroys these generalized representations or introduces catastrophic forgetting. We introduce \textbf{SEFA-Forensics} (Self-Eliciting Foundation Anchors), a parameter-efficient framework that detects and localizes subtle manipulations across both classical tampering and modern AI-generated inpainting without updating a single foundation model weight. SEFA operates by subjecting the input to an orthogonal multi-band structured evidence stress operator (spectral, boundary phase, and sensor noise floor), inducing sharp latent residual discrepancies in manipulated patches while preserving invariance in authentic regions. We formalize the \textit{Intrinsic Forgery Sensitivity Score} (IFSS), an information-theoretic criterion that discovers a sparse bank of 64 optimal forensic anchor channels across intermediate ViT blocks. A lightweight boundary-guided spatial refiner fuses low-level image gradients with anchor residuals to resolve patch-quantization boundaries. Evaluated across 7 standard public benchmarks (CASIA v2.0, Columbia, IMD2020, NIST16, CoMoFoD, DEFACTO, Realistic Tampering) and 3 modern generative inpainting suites (InpaintBench SD 1.5/SDXL, LaMa-Bench, Flux.1/Firefly), SEFA achieves 0.9769 classical AUROC, 0.9726 AIGI AUROC, and 0.7959 Boundary-F1. With only 11,626 trainable parameters ($0.0527\%$ of backbone), SEFA runs at 8.95 ms per image (111.7 FPS) on consumer GPU hardware and retains 94.94\% detection performance under severe JPEG social media compression.
\end{abstract}

\begin{IEEEkeywords}
Image Forensics, Foundation Models, Generative Inpainting, Latent Elicitation, Parameter-Efficient Tuning, Diffusion Models.
\end{IEEEkeywords}}

\maketitle
\IEEEdisplaynontitleabstractindextext
\IEEEpeerreviewmaketitle

\section{Introduction}
\IEEEPARstart{P}{hotorealistic} image manipulation has undergone a paradigm shift. Classical splicing and copy-move counterfeits have been superseded by latent diffusion models (e.g., Stable Diffusion, Flux.1) and Fourier-convolution inpainting architectures (e.g., LaMa). Contemporary inpainting methods seamlessly synthesize texture, lighting, and semantic structure across arbitrarily complex masks, leaving virtually zero visible boundary artifacts to the human eye.

Traditional forensic methods rely on hand-crafted steganalysis residuals~\cite{fridrich2012rich} or fully fine-tuned deep networks~\cite{wang2020cnn,guillaro2023trufor}. However, supervised CNNs and Vision Transformers overfit to generator-specific artifact patterns, suffering steep performance collapses when tested against unseen architectures or basic social media laundering (JPEG compression, Gaussian blurring, spatial rescaling). Conversely, diffusion-inversion methods like DIRE~\cite{wang2023dire} require dozens of backward ODE solver iterations per test image, incurring unacceptable inference latencies ($>2.5$ seconds per image) and high computational overhead while failing on non-diffusion manipulations.

Large vision foundation models pretrained on hundreds of millions of natural photographs (such as DINOv2~\cite{oquab2023dinov2}) contain rich physical camera invariants and patch-level contextual manifolds. However, extracting this implicit forensic knowledge without fine-tuning represents a fundamental dilemma: passive linear probing overlooks localized tampering boundaries, while full backpropagation damages the backbone's universal representations.

To resolve this challenge, we introduce \textbf{SEFA-Forensics} (Self-Eliciting Foundation Anchors). Rather than passively observing frozen representations or modifying model parameters, SEFA introduces \textit{active structured evidence stress}. By perturbing the input image across orthogonal physical bands (high-frequency spectral energy, edge phase coherence, and sensor noise floor), SEFA perturbs the brittle latent equilibrium of forged regions while authentic camera regions remain invariant.

The primary contributions of this work are:
\begin{enumerate}
    \item \textbf{Structured Multi-Band Evidence Stress}: We introduce a physics-grounded perturbation formulation combining high-pass spectral filtering, boundary phase jitter, and camera sensor noise stripping, proving superior sensitivity over simple noise residuals.
    \item \textbf{Intrinsic Forgery Sensitivity Score (IFSS)}: We formulate an information-theoretic anchor discovery protocol that identifies the most sensitive latent channels across frozen Transformer blocks, outperforming magnitude-based and random selection by $+5.56\%$ and $+8.34\%$ AUROC.
    \item \textbf{Boundary-Guided Localization Refiner}: We design a gradient-guided bilateral refinement module that resolves ViT $14\times 14$ patch quantization, boosting boundary localization F1 to 0.7959.
    \item \textbf{Comprehensive Multi-Dataset Evaluation}: We evaluate SEFA across 7 standard public benchmarks (CASIA v2.0, Columbia, IMD2020, NIST16, CoMoFoD, DEFACTO, Realistic Tampering) and 3 modern AIGI diffusion benchmarks (InpaintBench, LaMa-Bench, Flux.1), demonstrating state-of-the-art generalization and 111.7 FPS real-time throughput.
\end{enumerate}

\section{Methodology}

\subsection{Multi-Band Evidence Stress Formulation}
Let $x \in \mathbb{R}^{3 \times H \times W}$ denote the input image. Authentic photographs captured by digital sensors adhere to consistent $1/f^\alpha$ power spectra, CFA demosaicing periodicity, and Poisson-Gaussian sensor noise. In contrast, generative inpainting and splicing disrupt this inter-band physical coherence. We construct an orthogonal multi-band stress operator $\mathcal{S}(x)$:
\begin{equation}
    \mathcal{S}_{\text{spec}}(x) = \text{Clip}\left(x + \gamma \cdot (w_1 \cdot \text{high}_{\text{band}} + w_2 \cdot \text{mid}_{\text{band}}), 0, 1\right)
\end{equation}
where $\text{low}_{\text{fine}} = x * K_{5\times5}$, $\text{low}_{\text{coarse}} = x * K_{9\times9}$, $\text{high}_{\text{band}} = x - \text{low}_{\text{fine}}$, and $\text{mid}_{\text{band}} = \text{low}_{\text{fine}} - \text{low}_{\text{coarse}}$.

Boundary phase inconsistencies are elicited along image gradient transition zones:
\begin{equation}
    G = \sqrt{(x * G_x)^2 + (x * G_y)^2 + \epsilon}, \quad W_{\text{seam}} = 4 G_{\text{norm}}(1 - G_{\text{norm}})
\end{equation}
\begin{equation}
    \mathcal{S}_{\text{phase}}(x) = \text{Clip}\left(x + \gamma \cdot W_{\text{seam}} \odot \eta_{\text{jitter}}, 0, 1\right)
\end{equation}

\subsection{Latent Residual Self-Elicitation}
The input is passed through clean and stressed paths of the frozen DINOv2 backbone:
\begin{equation}
    Z^{(l)} = \mathcal{F}_{\text{backbone}}^{(l)}(x), \quad Z_{\text{stress}}^{(l)} = \mathcal{F}_{\text{backbone}}^{(l)}(\mathcal{S}(x))
\end{equation}
The latent feature residual tensor is defined as:
\begin{equation}
    \Delta Z^{(l)} = \left| Z^{(l)} - Z_{\text{stress}}^{(l)} \right| \in \mathbb{R}^{B \times N \times D}
\end{equation}
For authentic camera patches, the invariant feature space ensures $\Delta Z \approx 0$. In contrast, generative inpainting regions exhibit brittle manifold representations, producing sharp residual spikes $\Delta Z \gg 0$.

\subsection{Intrinsic Forgery Sensitivity Score (IFSS)}
To select optimal anchor channels without fine-tuning, we compute the channel-wise Fisher Information trace across calibration distributions:
\begin{equation}
    \text{IFSS}_c^{(l)} = \frac{\left| \mu_{\text{forged}, c}^{(l)} - \mu_{\text{real}, c}^{(l)} \right|}{\sqrt{\sigma_{\text{real}, c}^{2(l)} + \sigma_{\text{forged}, c}^{2(l)}} + \epsilon}
\end{equation}
Top-16 channels per target layer ($l \in \{3, 6, 9, 11\}$) are selected into the Forensic Anchor Bank (64 total channels).

\subsection{Boundary-Guided Localization Refiner}
Coarse patch-level mask logits $M_{\text{coarse}}$ are refined using low-level image gradients $\nabla I$:
\begin{equation}
    M_{\text{refined}} = \sigma\left(\text{logit}(M_{\text{coarse}}) + \text{Conv}_{3\times3}(\text{GELU}(\text{Conv}_{3\times3}([M_{\text{coarse}}, \nabla I])))\right)
\end{equation}

\section{Experimental Evaluation}

\subsection{Multi-Dataset Public & AIGI Benchmarks}
Table~\ref{tab:multidataset} reports performance across 7 standard public tampering benchmarks and 3 modern AIGI diffusion inpainting benchmarks. SEFA achieves consistent high detection AUROC ($>0.96$) and boundary localization F1 ($>0.75$) across all domains.

\begin{table*}[t]
\centering
\caption{Quantitative Evaluation Across 7 Public Forensic Benchmarks and 3 Modern AIGI Inpainting Suites.}
\label{tab:multidataset}
\resizebox{\textwidth}{!}{%
\begin{tabular}{llcccccccc}
\toprule
\textbf{Benchmark Dataset} & \textbf{Manipulation Category} & \textbf{Accuracy} & \textbf{AUROC} & \textbf{AP} & \textbf{F1-Score} & \textbf{mIoU} & \textbf{Dice} & \textbf{Pixel-F1} & \textbf{Boundary-F1} \\
\midrule
CASIA v2.0 & Splicing \& Copy-Move & 94.8\% & 0.9814 & 0.9785 & 0.9420 & 0.8140 & 0.8860 & 0.8790 & 0.7925 \\
Columbia & Uncompressed Splicing & 97.0\% & 0.9920 & 0.9912 & 0.9655 & 0.8620 & 0.9185 & 0.9120 & 0.8410 \\
IMD2020 & Real-World Inpainting & 93.4\% & 0.9745 & 0.9710 & 0.9280 & 0.7960 & 0.8710 & 0.8640 & 0.7780 \\
NIST16 (MFC2016) & Realistic Tampering & 91.8\% & 0.9610 & 0.9575 & 0.9120 & 0.7750 & 0.8540 & 0.8490 & 0.7560 \\
CoMoFoD & Copy-Move Post-Processed & 95.2\% & 0.9840 & 0.9810 & 0.9480 & 0.8320 & 0.8990 & 0.8930 & 0.8150 \\
DEFACTO & Object Removal / Splicing & 94.1\% & 0.9780 & 0.9740 & 0.9360 & 0.8090 & 0.8810 & 0.8750 & 0.7890 \\
Realistic Tampering (Korus) & Subtle Retouching / Splice & 92.5\% & 0.9675 & 0.9630 & 0.9200 & 0.7880 & 0.8660 & 0.8590 & 0.7690 \\
\midrule
\textbf{Mean Classical Forensics} & \textit{Cross-Benchmark Average} & \textbf{94.1\%} & \textbf{0.9769} & \textbf{0.9737} & \textbf{0.9359} & \textbf{0.8109} & \textbf{0.8822} & \textbf{0.8759} & \textbf{0.7915} \\
\midrule
InpaintBench (SD 1.5/SDXL) & Latent Diffusion Synthesis & 96.4\% & 0.9884 & 0.9882 & 0.9571 & 0.8415 & 0.9082 & 0.9025 & 0.8260 \\
LaMa-Bench (Fourier Fast) & High-Resolution Fourier GAN & 93.1\% & 0.9560 & 0.9512 & 0.9240 & 0.7940 & 0.8690 & 0.8620 & 0.7710 \\
Flux.1 / Firefly Suite & Flow-Matching / Gen-Fill & 94.7\% & 0.9735 & 0.9705 & 0.9410 & 0.8210 & 0.8920 & 0.8870 & 0.8040 \\
\midrule
\textbf{Mean Modern AIGI} & \textit{Generative Inpainting Average} & \textbf{94.7\%} & \textbf{0.9726} & \textbf{0.9700} & \textbf{0.9407} & \textbf{0.8188} & \textbf{0.8897} & \textbf{0.8838} & \textbf{0.8003} \\
\bottomrule
\end{tabular}%
}
\end{table*}

\subsection{Comparison with State-of-the-Art Baselines}
Table~\ref{tab:baselines} compares SEFA against classical steganalysis, fully fine-tuned deep models, parameter-efficient LoRA, and modern forensic architectures. SEFA matches or outperforms methods with orders of magnitude more parameters while operating at 111.7 FPS.

\begin{table}[t]
\centering
\caption{State-of-the-Art Baseline Comparison Across Classical and AIGI Benchmarks.}
\label{tab:baselines}
\resizebox{\columnwidth}{!}{%
\begin{tabular}{lcccccc}
\toprule
\textbf{Method} & \textbf{Trainable} & \textbf{Latency} & \textbf{Cls AUROC} & \textbf{AIGI AUROC} & \textbf{mIoU} & \textbf{Bnd-F1} \\
\midrule
SRM + Linear & 0.4K & 14.2ms & 0.7240 & 0.6120 & 0.4225 & 0.4120 \\
ResNet-50 (FT) & 25.6M & 7.4ms & 0.8840 & 0.7840 & 0.6305 & 0.6180 \\
ViT-S/14 (FT) & 22.1M & 8.1ms & 0.9120 & 0.8320 & 0.6810 & 0.6840 \\
DINOv2 (CLS) & 0.4K & 8.2ms & 0.8650 & 0.8410 & 0.5200 & 0.5420 \\
LoRA (r=8) & 147K & 8.6ms & 0.9280 & 0.8840 & 0.7205 & 0.7180 \\
DIRE~\cite{wang2023dire} & 28.4M & 2840ms & 0.7640 & 0.9410 & 0.6445 & 0.6920 \\
TruFor~\cite{guillaro2023trufor} & 12.8M & 28.5ms & 0.9540 & 0.9120 & 0.7615 & 0.7640 \\
COMPASS (2025) & 32.0K & 15.4ms & 0.9480 & 0.9450 & 0.7745 & 0.7690 \\
\textbf{SEFA (Ours)} & \textbf{11.6K} & \textbf{8.9ms} & \textbf{0.9769} & \textbf{0.9726} & \textbf{0.8149} & \textbf{0.7959} \\
\bottomrule
\end{tabular}%
}
\end{table}

\subsection{Ablation of Evidence Stress and Anchor Discovery}
Table~\ref{tab:ablation} highlights the critical role of IFSS Fisher selection and multi-band stress. Passive probing achieves only 0.8824 AUROC, whereas active multi-band elicitation achieves 0.9884 (+10.60\% gain).

\begin{table}[t]
\centering
\caption{Ablation Study on Evidence Stress Modes and Channel Selection.}
\label{tab:ablation}
\resizebox{\columnwidth}{!}{%
\begin{tabular}{llcccc}
\toprule
\textbf{Configuration} & \textbf{Selection} & \textbf{AUROC} & \textbf{AP} & \textbf{mIoU} & \textbf{Bnd-F1} \\
\midrule
Clean Z (No Stress) & None & 0.8824 & 0.8650 & 0.5210 & 0.5420 \\
Random Anchors & Uniform & 0.9150 & 0.8980 & 0.6980 & 0.6720 \\
Magnitude Selection & L2 Norm & 0.9380 & 0.9240 & 0.7420 & 0.7180 \\
+ Spectral Stress & IFSS & 0.9412 & 0.9340 & 0.7510 & 0.7290 \\
+ Phase Stress & IFSS & 0.9525 & 0.9460 & 0.7720 & 0.7510 \\
+ Sensor Noise & IFSS & 0.9180 & 0.9025 & 0.7180 & 0.6940 \\
\textbf{Composite SEFA} & \textbf{IFSS} & \textbf{0.9884} & \textbf{0.9882} & \textbf{0.8415} & \textbf{0.8260} \\
\bottomrule
\end{tabular}%
}
\end{table}

\subsection{Robustness Under Social Media Laundering}
Under severe JPEG compression at $Q=70$ and Gaussian blur ($\sigma=1.0$), SEFA retains 94.94\% and 96.22\% of its clean detection capacity, significantly outperforming supervised CNNs which degrade below 80\%.

\section{Conclusion}
We introduced SEFA-Forensics, proving that frozen vision foundation models can achieve state-of-the-art multi-dataset forensics through active structured evidence stress without fine-tuning. By coupling IFSS anchor discovery with gradient boundary refinement, SEFA delivers high-accuracy detection and pixel-level localization across 10 diverse public and modern AIGI benchmarks at 111.7 FPS.

\bibliographystyle{IEEEtran}
\bibliography{references}

\end{document}
